% RecSys Odyssey repository-backed lecture note.
% Add figures with: \coursefigure{figures/name.svg}{Caption}
\section{The Transformer block}

\begin{abstract}
Attention mixes positions; the feed-forward network transforms each position.
\end{abstract}

\subsection{Concept}
A Transformer block combines multi-head attention with a position-wise feed-forward network. Residual connections preserve information and stabilize gradients; normalization keeps activations well behaved. Stacking blocks expands the depth of interaction the model can express. The final representation of the last token, a special token or a pooled sequence becomes the user state used for retrieval or ranking.

\begin{equation*}
H′ = Norm(H + MHA(H))      H″ = Norm(H′ + FFN(H′))
\end{equation*}

\subsection{Key terms}
\begin{itemize}
\item \textbf{Residual connection.} A skip path that adds a block’s input to its output.
\item \textbf{Feed-forward network.} A per-position nonlinear transformation after attention.
\item \textbf{Layer normalization.} Feature-wise normalization that stabilizes optimization.
\end{itemize}
